% ============================================================================
% CHAPTER 1: Projektbeschreibung
% ============================================================================
\renewcommand{\chaptertag}{Bildfilter}
\chapter{Projektbeschreibung}

\vspace{-5pt}
{\color{maincolor}\rule{\textwidth}{0.4pt}}
\vspace{4pt}

{\color{maincolor}\textbf{Lernziele:}}
\begin{itemize}[leftmargin=*]
  \item Du verstehst, wie ein Computer ein Bild speichert.
  \item Du kannst eigene Bildfilter als Python-Funktionen schreiben und an
    echten Fotos ausprobieren.
  \item Du gehst Schritt f\"ur Schritt von einfachen Pixel-\"Anderungen
    bis zu Filtern, die mehrere Nachbarpixel zusammenrechnen.
\end{itemize}
\vspace{-2pt}
{\color{maincolor}\rule{\textwidth}{0.4pt}}
\vspace{8pt}

Jedes Mal, wenn du auf Instagram einen Filter anwendest oder in der
Galerie-App ein Foto heller machst, l\"auft im Hintergrund ein
Bildbearbeitungs-Algorithmus. In diesem Projekt programmierst du solche
Filter selbst und siehst sofort, was dein Code mit echten Bildern macht.

% ============================================================================
\section{So funktioniert das Projekt}

\begin{enumerate}
  \item Du arbeitest an deinem eigenen Python-Programm. Du schreibst Filter
    als Funktionen.
  \item Auf einer Website f\"uhrst du ein \textit{Stempelheft}. F\"ur jede
    Aufgabe, die du schaffst, bekommst du einen Stempel.
  \item Mit jedem Stempel wird ein St\"uck eines paint-by-numbers-Bildes auf
    deinem Heft sichtbar. Wer alle Stempel hat, sieht das fertige Bild.
\end{enumerate}

\begin{definitionbox}{Bildfilter}
Ein \term{Bildfilter} ist ein Algorithmus, der ein Bild als Eingabe nimmt
und ein neues Bild zur\"uckgibt. Dabei wird jeder Pixel nach einer
bestimmten Regel ver\"andert.
\end{definitionbox}

% ============================================================================
\section{Spielregeln}

\begin{itemize}
  \item Du arbeitest mit \textbf{verschachtelten Listen}, wie du sie aus
    dem Algorithmen-Skript kennst:
\begin{lstlisting}
bild[zeile][spalte]  # liefert einen Pixel (R, G, B)
\end{lstlisting}
  \item Zum Laden und Speichern von Bilddateien verwendest du die
    Bibliothek \texttt{Pillow}. Setze die folgenden zwei Funktionen
    einmal an den Anfang deiner Datei. Sie laden ein Bild aus dem
    aktuellen Ordner in eine verschachtelte Liste und speichern eine
    verschachtelte Liste wieder als Bild:
\begin{lstlisting}
from PIL import Image

def lade_bild(pfad):
    img = Image.open(pfad).convert("RGB")
    breite, hoehe = img.size
    bild = []
    for y in range(hoehe):
        zeile = []
        for x in range(breite):
            zeile.append(img.getpixel((x, y)))
        bild.append(zeile)
    return bild

def speichere_bild(bild, pfad):
    hoehe = len(bild)
    breite = len(bild[0])
    img = Image.new("RGB", (breite, hoehe))
    for y in range(hoehe):
        for x in range(breite):
            img.putpixel((x, y), bild[y][x])
    img.save(pfad)
\end{lstlisting}
    Aufrufen kannst du sie so:
\begin{lstlisting}
bild = lade_bild("foto.jpg")
speichere_bild(bild, "ergebnis.png")
\end{lstlisting}
  \item Alles andere (jede Filterlogik) schreibst du selbst.
\end{itemize}

\begin{beispiel}{Was ist ein Pixel?}
Ein Pixel wird in Python als Tupel aus drei Zahlen gespeichert, jeweils
zwischen 0 und 255:

\begin{lstlisting}
schwarz = (0, 0, 0)
weiss   = (255, 255, 255)
rot     = (255, 0, 0)
gruen   = (0, 255, 0)
blau    = (0, 0, 255)
gelb    = (255, 255, 0)   # Rot + Gruen
violett = (128, 0, 128)
\end{lstlisting}

Auf einen einzelnen Pixel greifst du so zu:
\begin{lstlisting}
pixel = bild[10][20]   # Pixel in Zeile 10, Spalte 20
r, g, b = pixel        # in seine drei Werte zerlegen
\end{lstlisting}
\end{beispiel}

% ============================================================================
\section{Abgabe}

Du gibst am Ende drei Sachen ab:

\begin{enumerate}
  \item Deine \textbf{Python-Datei(en)} mit allen Filtern.
  \item Drei \textbf{Vorher/Nachher-Bildpaare}: ein Originalfoto und
    daneben jeweils das Ergebnis eines deiner Filter.
  \item Einen \textbf{Bericht} (2 bis 3 A4-Seiten): Was hast du
    gemacht? Welcher Filter hat dich \"uberrascht? Wo bist du steckengeblieben?
\end{enumerate}

Die Abgabe l\"auft \"uber die normale Plattform der Klasse. Wie genau,
besprechen wir gemeinsam.

% ============================================================================
\section{Zeitrahmen}

Das Projekt l\"auft \"uber rund sieben Doppellektionen. Das n\"achste
Kapitel zeigt dir, wie bewertet wird; danach folgt der \textit{Lernpfad}
mit den eigentlichen Aufgaben.

\closechapter
